\documentclass[12pt,a4paper]{article}
\usepackage[utf8]{inputenc}
\usepackage[margin=1in]{geometry}
\usepackage{graphicx}
\usepackage{hyperref}
\usepackage{amsmath}
\usepackage{amssymb}
\usepackage{booktabs}
\usepackage{enumitem}
\usepackage{float}
\usepackage{caption}
\usepackage{listings}
\usepackage{xcolor}
\usepackage{titlesec}
\usepackage{fancyhdr}
\usepackage{multirow}
\usepackage{array}
\usepackage{tabularx}

% Colors
\definecolor{codeblue}{RGB}{30, 100, 200}
\definecolor{codegray}{RGB}{128,128,128}
\definecolor{codegreen}{RGB}{40,140,40}
\definecolor{codebg}{RGB}{245,245,250}
\definecolor{accent}{RGB}{46,134,171}
\definecolor{darkblue}{RGB}{29,53,87}

% Code listing style
\lstset{
  backgroundcolor=\color{codebg},
  basicstyle=\ttfamily\small,
  breaklines=true,
  frame=single,
  rulecolor=\color{accent},
  keywordstyle=\color{codeblue}\bfseries,
  commentstyle=\color{codegreen}\itshape,
  stringstyle=\color{red!60!black},
  showstringspaces=false,
  numbers=left,
  numberstyle=\tiny\color{codegray},
  numbersep=8pt,
  tabsize=4,
  language=Python,
  captionpos=b,
  belowskip=1em,
  aboveskip=1em
}

\hypersetup{
    colorlinks=true,
    linkcolor=darkblue,
    urlcolor=accent,
    citecolor=darkblue
}

\titleformat{\section}{\large\bfseries\color{darkblue}}{\thesection}{1em}{}[\titlerule]
\titleformat{\subsection}{\normalsize\bfseries\color{accent}}{\thesubsection}{1em}{}

\pagestyle{fancy}
\fancyhf{}
\rhead{\small ECSCI24305 --- Deep Learning}
\lhead{\small Rishi Raval --- 1AUA23BCS155}
\cfoot{\thepage}
\renewcommand{\headrulewidth}{0.4pt}

\begin{document}

%── Title Page ────────────────────────────────────────────────────────────────
\begin{titlepage}
\centering
\vspace*{1.5cm}
{\Large \textbf{Adani University}\\[0.3cm]
\large School of Engineering and Technology}\\[1cm]
\rule{\linewidth}{1pt}\\[0.5cm]
{\LARGE \textbf{Project-Based Learning Report}}\\[0.3cm]
{\large Phase 3 \& Phase 4 Submission}\\[0.3cm]
\rule{\linewidth}{1pt}\\[1cm]

{\Large \textbf{Multimodal Volatility \& Direction Forecasting\\
Using Deep Learning}}\\[0.5cm]

{\large \textit{Predicting stock market volatility through fusion of\\
price patterns, graph relationships, SEC filings, and macroeconomic data}}\\[1.5cm]

\begin{tabular}{ll}
\textbf{Student Name}   & Rishi Raval \\[0.3cm]
\textbf{Roll Number}    & 1AUA23BCS155 \\[0.3cm]
\textbf{Programme}      & B.Tech CSE (AI-ML), Semester VI \\[0.3cm]
\textbf{Course}         & ECSCI24305 --- Deep Learning: Principles and Practices \\[0.3cm]
\textbf{Submission}     & Phase 3 (Preprocessing) \& Phase 4 (Model Design) \\[0.3cm]
\textbf{Date}           & March 2026 \\[0.3cm]
\textbf{GitHub}         & \url{https://github.com/imrishi007/financial-document-analysis}
\end{tabular}

\vfill
\end{titlepage}

\tableofcontents
\newpage

%══════════════════════════════════════════════════════════════════════════════
\section{Project Introduction \& Motivation}
%══════════════════════════════════════════════════════════════════════════════

\subsection{Problem Statement}

Financial markets generate enormous volumes of heterogeneous data every trading day --- price time series, corporate filings, macroeconomic indicators, and inter-company relationships. This project builds an end-to-end multimodal deep learning system that fuses four distinct data modalities to answer two questions simultaneously:

\begin{enumerate}
    \item \textbf{Volatility Forecasting (Primary):} How wildly will a stock's price move over the next 20 trading days? Formally, predict the 20-day annualized realized volatility $\sigma_t$ defined as:
    \begin{equation}
        \sigma_t = \sqrt{\frac{252}{20} \sum_{i=1}^{20} r_{t+i}^2}
    \end{equation}
    where $r_{t+i} = \log(P_{t+i}/P_{t+i-1})$ is the daily log return.

    \item \textbf{Direction Forecasting (Secondary):} Will this stock's price be higher or lower in 60 trading days?
\end{enumerate}

\subsection{Motivation}

\textbf{Why volatility?} Volatility is the cornerstone of financial risk management. It determines option prices (Black-Scholes), portfolio Value-at-Risk, and margin requirements. Crucially, as Nobel laureate Robert Engle demonstrated in his foundational 1982 ARCH paper, \textit{volatility is persistent} --- turbulent markets tend to stay turbulent, and calm markets tend to stay calm. This autocorrelation makes volatility more predictable than direction, which is near-random on efficient markets (Efficient Market Hypothesis).

\textbf{Why multimodal?} No single data source tells the full story. Price patterns capture short-term momentum. SEC filings encode long-term fundamental outlook. Cross-stock graph relationships capture sector-wide dynamics. Macroeconomic indicators set the market regime. Our hypothesis is that fusing all four provides richer representations than any single modality.

\textbf{Why deep learning?} Traditional econometric models (GARCH, HAR-RV) capture volatility clustering through linear autoregression. Deep learning can capture nonlinear interactions between modalities --- for example, how a negative SEC filing interacts with a deteriorating macroeconomic environment to amplify volatility.

\subsection{Project Universe}

\begin{table}[H]
\centering
\caption{Stock universe (30 companies across 6 sectors)}
\begin{tabular}{ll}
\toprule
\textbf{Sector} & \textbf{Tickers} \\
\midrule
Core Tech (original)    & AAPL, MSFT, GOOGL, AMZN, META, NVDA, AMD, INTC, TSLA, ORCL \\
Semiconductors          & QCOM, TXN, AVGO, MRVL, KLAC \\
Cloud / Software        & CRM, ADBE, NOW, SNOW, DDOG \\
Internet / Consumer     & NFLX, UBER, PYPL, SNAP \\
Hardware / Infra        & DELL, AMAT, LRCX \\
Diversified Tech        & IBM, CSCO, HPE \\
\bottomrule
\end{tabular}
\end{table}

%══════════════════════════════════════════════════════════════════════════════
\section{Phase 3: Data Preprocessing \& Analysis}
%══════════════════════════════════════════════════════════════════════════════

\subsection{Dataset Overview}

The project uses four independent data sources, each requiring different preprocessing pipelines:

\begin{table}[H]
\centering
\caption{Summary of all data modalities}
\begin{tabular}{llll}
\toprule
\textbf{Modality} & \textbf{Source} & \textbf{Coverage} & \textbf{Samples} \\
\midrule
Price (OHLCV)       & Yahoo Finance    & 30 stocks, 2016--2026 & 73,793 windows \\
SEC 10-K Filings    & SEC EDGAR        & 227 annual filings    & 73.7\% of dates \\
Macroeconomic       & Yahoo Finance    & 12 indicators, daily  & 99.6\% of dates \\
Graph (Knowledge)   & Hand-curated     & 30 nodes, 200 edges   & 99.3\% of dates \\
\bottomrule
\end{tabular}
\end{table}

\subsection{Temporal Train/Validation/Test Split}

\textbf{Critical design decision:} All splits are strictly chronological. Random splitting is fundamentally incorrect for time series because it would allow a model trained on 2020 data to be tested on 2019 data, having already ``seen the future.'' This is data leakage.

\begin{table}[H]
\centering
\caption{Temporal data splits}
\begin{tabular}{lllr}
\toprule
\textbf{Split} & \textbf{Date Range} & \textbf{Samples} & \textbf{Percentage} \\
\midrule
Train      & 2016-01-01 -- 2022-12-31 & 48,532 & 65.8\% \\
Validation & 2023-01-01 -- 2023-12-31 &  8,057 & 10.9\% \\
Test       & 2024-01-01 -- 2025-12-31 & 17,204 & 23.3\% \\
\midrule
\textbf{Total} & & \textbf{73,793} & \textbf{100\%} \\
\bottomrule
\end{tabular}
\end{table}

\begin{figure}[H]
\centering
\includegraphics[width=0.65\textwidth]{figures/dataset_split.png}
\caption{Chronological train/validation/test split distribution}
\end{figure}

The split implementation uses date-based cutoffs, not ratio-based indexing:

\begin{lstlisting}[caption={Temporal split logic (src/data/preprocessing.py)}]
@dataclass
class SplitConfig:
    """Configuration for chronological train / val / test split."""
    train_end: str = "2022-12-31"
    val_end: str = "2023-12-31"
    # Everything after val_end is test

def create_time_splits(dates, cfg=None):
    """Return boolean index arrays for train / val / test
    based on date cutoffs."""
    if cfg is None:
        cfg = SplitConfig()
    dt = pd.to_datetime(dates)
    train_end = pd.Timestamp(cfg.train_end)
    val_end = pd.Timestamp(cfg.val_end)

    train_mask = (dt <= train_end).values
    val_mask = ((dt > train_end) & (dt <= val_end)).values
    test_mask = (dt > val_end).values
    return {"train": train_mask, "val": val_mask,
            "test": test_mask}
\end{lstlisting}

\subsection{Price Data Preprocessing}

\subsubsection{Feature Engineering (21 Features)}

Raw OHLCV data is transformed into 21 engineered features per timestep. These fall into three groups:

\textbf{Group 1: Standard Technical Features (10)}
\begin{enumerate}[noitemsep]
    \item Normalized close price
    \item Normalized volume
    \item 1-day return: \texttt{ret\_1}
    \item 5-day return: \texttt{ret\_5}
    \item 10-day return: \texttt{ret\_10}
    \item 10-day rolling volatility: \texttt{volatility\_10}
    \item Volume change: \texttt{volume\_change\_1}
    \item 5-day moving average: \texttt{ma\_5}
    \item 20-day moving average: \texttt{ma\_20}
    \item MA ratio: \texttt{ma\_ratio\_5\_20}
\end{enumerate}

\textbf{Group 2: Volatility-Specific Features (8)}
\begin{enumerate}[noitemsep,resume]
    \item Parkinson volatility (high/low range estimator, 10-day)
    \item Garman-Klass volatility (OHLC estimator, 10-day)
    \item 5-day realized volatility: \texttt{rv\_5d}
    \item Volatility ratio: $\sigma_{10d}/\sigma_{20d}$ (regime indicator)
    \item Squared return: $r_t^2$ (instantaneous variance proxy)
    \item Absolute return: $|r_t|$ (robust volatility estimator)
    \item Jump indicator: $\mathbb{1}[|r_t| > 3\sigma_{20d}]$
    \item EWMA variance (span=10)
\end{enumerate}

\textbf{Group 3: HAR-RV Autoregressive Features (3) --- Key Innovation}
\begin{enumerate}[noitemsep,resume]
    \item \texttt{rv\_lag1d}: Yesterday's realized volatility (daily component)
    \item \texttt{rv\_lag5d}: 5-day rolling mean RV (weekly component)
    \item \texttt{rv\_lag22d}: 22-day rolling mean RV (monthly component)
\end{enumerate}

These three features directly encode the Heterogeneous Autoregressive (HAR) structure of Corsi (2009), which models volatility persistence at three timescales. All three use \texttt{shift(1)} to prevent look-ahead bias.

\begin{lstlisting}[caption={Price feature engineering --- HAR-RV features (src/features/price\_features.py)}]
# --- Phase 13: 3 HAR-RV autoregressive features ---
# HAR-RV model (Corsi 2009):
#   RV_t = b0 + b_d*RV_{t-1} + b_w*RV^w_{t-5} + b_m*RV^m_{t-22}
log_ret = np.log(
    frame["close"]
    / frame["close"].groupby(frame["ticker"]).shift(1)
)
# 20-day rolling realized vol, annualized
rv_base = (
    log_ret.groupby(frame["ticker"])
    .rolling(20).std()
    .reset_index(level=0, drop=True) * np.sqrt(252)
)

# HAR-RV Feature 1: daily lag (t-1)
rv_lag1d = rv_base.groupby(frame["ticker"]).shift(1) \
    .reset_index(level=0, drop=True)
# HAR-RV Feature 2: weekly average (t-1 to t-5)
rv_lag5d = (
    rv_base.groupby(frame["ticker"]).rolling(5).mean()
    .reset_index(level=0, drop=True)
    .groupby(frame["ticker"]).shift(1)
    .reset_index(level=0, drop=True)
)
# HAR-RV Feature 3: monthly average (t-1 to t-22)
rv_lag22d = (
    rv_base.groupby(frame["ticker"]).rolling(22).mean()
    .reset_index(level=0, drop=True)
    .groupby(frame["ticker"]).shift(1)
    .reset_index(level=0, drop=True)
)
\end{lstlisting}

\subsubsection{Z-Score Normalization}

All features are normalized using statistics computed on the training set only:

\begin{equation}
    x_{\text{norm}} = \frac{x - \mu_{\text{train}}}{\sigma_{\text{train}} + \epsilon}, \quad \epsilon = 10^{-8}
\end{equation}

\begin{lstlisting}[caption={Feature normalization (src/data/preprocessing.py)}]
@dataclass
class FeatureScaler:
    """Stores mean / std computed on training data
    for z-score normalization."""
    mean: pd.Series
    std: pd.Series

    def transform(self, df, cols):
        out = df.copy()
        for c in cols:
            out[c] = (out[c] - self.mean[c]) \
                      / (self.std[c] + 1e-8)
        return out

def fit_scaler(df, cols):
    """Fit a z-score scaler on the training set only."""
    return FeatureScaler(
        mean=df[cols].mean(), std=df[cols].std()
    )
\end{lstlisting}

\textbf{Why training set only?} Using validation or test statistics would leak future distribution information into the model.

\subsubsection{Sliding Window Dataset}

Each training sample is a 60-day window of the 21 features, producing a tensor of shape $[60, 21]$:

\begin{lstlisting}[caption={PriceWindowDataset (src/data/price\_dataset.py)}]
class PriceWindowDataset(Dataset):
    def __init__(self, price_df, targets_df,
                 vol_df=None, surprise_df=None,
                 window_size=60, feature_cols=None):
        self.window_size = window_size
        self.feature_cols = feature_cols \
            or ENGINEERED_FEATURES  # 21 features

        # Merge price features with direction targets
        merged = price_df.merge(
            targets_df, on=["ticker", "date"],
            how="inner")
        # Optionally merge volatility and surprise
        if vol_df is not None:
            merged = merged.merge(vol_df, ...)
        # Build sliding windows per ticker
        self._samples = []
        for ticker, group in merged.groupby("ticker"):
            for i in range(window_size, len(group)):
                window = group.iloc[
                    i - window_size : i]
                target_row = group.iloc[i]
                sample = {
                    "features": window[
                        self.feature_cols].values,
                    "direction_60d": int(
                        target_row["direction_60d_id"]),
                    "volatility": float(target_row[
                        "realized_vol_20d_annualized"]),
                    "ticker": str(ticker),
                    "date": str(target_row["date"]),
                }
                self._samples.append(sample)
\end{lstlisting}

\subsection{SEC 10-K Filing Preprocessing}

Annual SEC filings are long-form documents (often 100+ pages). FinBERT has a maximum context length of 512 tokens, so each filing is chunked:

\begin{itemize}[noitemsep]
    \item \textbf{Max sequence length:} 512 tokens per chunk (510 + [CLS] + [SEP])
    \item \textbf{Stride (overlap):} 128 tokens (preserves context at chunk boundaries)
    \item \textbf{Max chunks per document:} 64
    \item \textbf{Total filings processed:} 227 (across 30 stocks, 2016--2025)
\end{itemize}

\begin{lstlisting}[caption={Text chunking (src/data/preprocessing.py)}]
def chunk_text(text, tokenizer, max_length=510,
               stride=128):
    """Split long text into overlapping token chunks
    suitable for BERT."""
    encoded = tokenizer(text,
        add_special_tokens=False,
        return_attention_mask=False,
        return_token_type_ids=False)
    all_ids = encoded["input_ids"]

    cls_id = tokenizer.cls_token_id
    sep_id = tokenizer.sep_token_id

    chunks = []
    start = 0
    while start < len(all_ids):
        end = start + max_length
        chunk_ids = all_ids[start:end]
        # Add [CLS] and [SEP]
        input_ids = [cls_id] + chunk_ids + [sep_id]
        attention_mask = [1] * len(input_ids)
        # Pad to max_length + 2
        pad_len = (max_length + 2) - len(input_ids)
        if pad_len > 0:
            input_ids += [tokenizer.pad_token_id] \
                          * pad_len
            attention_mask += [0] * pad_len
        chunks.append({"input_ids": input_ids,
                       "attention_mask": attention_mask})
        if end >= len(all_ids):
            break
        start += max_length - stride
    return chunks
\end{lstlisting}

\textbf{Filing-to-date alignment:} Each 10-K filing from year $Y$ is mapped to all prediction dates in year $Y+1$, since annual reports become public early in the following year.

\subsection{Macroeconomic Data Preprocessing}

12 macroeconomic indicators are downloaded daily from Yahoo Finance and preprocessed:

\begin{table}[H]
\centering
\caption{Macroeconomic features}
\begin{tabular}{lll}
\toprule
\textbf{Feature} & \textbf{Source Ticker} & \textbf{Interpretation} \\
\midrule
VIX level        & \^{}VIX  & Market fear index \\
VIX daily change & \^{}VIX  & Direction of fear \\
Long-term yields & TLT      & 20-year Treasury ETF \\
Medium-term yields & IEF    & 7--10 year Treasury ETF \\
Yield curve spread & TLT-IEF & Recession indicator \\
Broad market momentum & SPY & S\&P 500 direction \\
Tech sector momentum  & QQQ & NASDAQ direction \\
Short-term rates  & SHY     & 1--3 year Treasury \\
Gold momentum    & GLD      & Safe haven demand \\
High-yield credit & HYG    & Credit stress \\
Investment-grade credit & LQD & Credit quality \\
Market breadth   & Computed & Cross-stock co-movement \\
\bottomrule
\end{tabular}
\end{table}

\textbf{Critical leakage prevention:} All macro features are lagged by 1 day (\texttt{shift(1)}). Today's VIX close is used to predict tomorrow's volatility, not today's.

\subsection{Graph Construction}

A static knowledge graph encodes business relationships between all 30 stocks:

\begin{itemize}[noitemsep]
    \item \textbf{Nodes:} 30 stocks
    \item \textbf{Directed edges:} 70 (based on business relationships)
    \item \textbf{Total edges:} 200 (bidirectional + self-loops)
    \item \textbf{Relation types:} supplier\_to, platform\_competitor, ad\_competitor, cloud\_competitor, consumer\_attention\_overlap, enterprise\_competitor, index\_correlation
\end{itemize}

\textbf{Why static, not dynamic?} Dynamic correlation-based graphs (rolling Pearson correlations) were tested in Phase 9 and performed worse (AUC dropped from 0.531 to 0.498). Curated business relationships are more stable and semantically meaningful than noisy rolling correlations.

\subsection{Modality Availability \& Masking}

Not all modalities are available for every sample. The fusion model handles missing modalities through availability masks:

\begin{figure}[H]
\centering
\includegraphics[width=0.75\textwidth]{figures/modality_coverage.png}
\caption{Modality coverage across all 73,793 samples}
\end{figure}

\subsection{Exploratory Data Analysis}

\subsubsection{Volatility Distribution}

\begin{figure}[H]
\centering
\includegraphics[width=\textwidth]{figures/volatility_analysis.png}
\caption{Left: Volatility distribution by stock (TSLA has highest, MSFT lowest). Right: Volatility clustering --- turbulent periods persist, validating Engle's ARCH insight.}
\end{figure}

Key observations from EDA:
\begin{itemize}[noitemsep]
    \item TSLA exhibits the highest mean volatility ($\sim$65\% annualized) due to idiosyncratic factors
    \item MSFT and AAPL cluster around 25--30\% annualized volatility
    \item Volatility clustering is clearly visible --- periods of high volatility persist for weeks
    \item The distribution is right-skewed (log-normal), consistent with financial literature
\end{itemize}

\subsubsection{Data Leakage Audit}

During development, two critical data leakages were identified and corrected:

\begin{enumerate}
    \item \textbf{Surprise Feature Leakage:} The original surprise feature vector $[\texttt{is\_beat}, \texttt{is\_miss}, \texttt{recency}]$ directly encoded the BEAT/MISS earnings label, producing a perfect AUC of 1.0. Fix: removed \texttt{is\_beat} and \texttt{is\_miss}, retained only recency.

    \item \textbf{Cross-Stock Feature Leakage:} Sector momentum features accidentally used \texttt{direction\_5d\_return} (forward-looking) from peer stocks, producing an inflated AUC of 0.939. Fix: replaced with \texttt{log\_return} (historical) and computed rolling past sums.
\end{enumerate}

Both leakages were caught by the same instinct: \textit{when a result is too good to be true, audit the input features immediately.}

%══════════════════════════════════════════════════════════════════════════════
% PHASE 4 continues in next input section
%══════════════════════════════════════════════════════════════════════════════
\section{Phase 4: Model Design \& Implementation}
%══════════════════════════════════════════════════════════════════════════════

\subsection{Architecture Overview}

The system consists of four modality-specific encoders whose outputs are fused by a learned gating mechanism. Two task heads produce the final predictions.

\begin{figure}[H]
\centering
\includegraphics[width=\textwidth]{figures/architecture.png}
\caption{Complete multimodal fusion architecture. The HAR-RV skip connection (red arrow) bypasses the CNN-BiLSTM bottleneck to directly feed autoregressive features into the fusion trunk.}
\end{figure}

\subsection{Model 1: CNN-BiLSTM Price Encoder}

\subsubsection{Design Rationale}

\begin{itemize}[noitemsep]
    \item \textbf{1D-CNN:} Detects local patterns in 3-day windows (e.g., momentum, reversals). Uses 64 filters to learn 64 different pattern types simultaneously.
    \item \textbf{BiLSTM:} Captures how CNN patterns evolve over the full 60-day window. Bidirectional processing gives context from both past and future within the window.
    \item \textbf{Why both?} CNN handles local structure; BiLSTM handles temporal evolution. Together they capture patterns at multiple timescales.
\end{itemize}

\begin{lstlisting}[caption={CNN-BiLSTM price encoder (src/models/price\_model.py)}]
class PriceDirectionModel(nn.Module):
    def __init__(self, num_features, conv_channels=64,
                 lstm_hidden=128):
        super().__init__()
        # Local pattern extraction: two Conv1d layers
        self.conv = nn.Sequential(
            nn.Conv1d(num_features, conv_channels,
                      kernel_size=3, padding=1),
            nn.ReLU(),
            nn.BatchNorm1d(conv_channels),
            nn.Conv1d(conv_channels, conv_channels,
                      kernel_size=3, padding=1),
            nn.ReLU(),
            nn.BatchNorm1d(conv_channels),
        )
        # Temporal evolution: 2-layer BiLSTM
        self.lstm = nn.LSTM(
            input_size=conv_channels,
            hidden_size=lstm_hidden,
            num_layers=2,
            batch_first=True,
            dropout=0.2,
            bidirectional=True,  # 128*2 = 256-d output
        )
        # Classification head
        self.head = nn.Sequential(
            nn.Linear(lstm_hidden * 2, 128),
            nn.ReLU(),
            nn.Dropout(0.2),
            nn.Linear(128, 2),
        )

    def forward(self, x):
        # x: [B, 60, 21]
        x = x.transpose(1, 2)   # [B, 21, 60] for Conv1d
        x = self.conv(x)        # [B, 64, 60]
        x = x.transpose(1, 2)   # [B, 60, 64] for LSTM
        x, _ = self.lstm(x)     # [B, 60, 256]
        x = x[:, -1, :]         # [B, 256] last timestep
        return self.head(x)     # [B, 2]
\end{lstlisting}

\textbf{Output:} 256-dimensional embedding per stock per day.

\subsection{Model 2: Graph Attention Network (GAT)}

\subsubsection{Design Rationale}

Stocks do not move in isolation. NVIDIA's supply relationships with Apple and Microsoft, or Google's competition with Meta in digital advertising, create cross-stock dependencies that a stock-independent model cannot capture. The GAT propagates information across the 30-stock graph, allowing each stock's representation to incorporate signals from its business neighbors.

\textbf{Key constraint:} Implemented in pure PyTorch without PyTorch Geometric, as required by the project specification.

\subsubsection{GAT Attention Mechanism}

Following Veli\v{c}kovi\'c et al.\ (2018), each GAT layer computes:

\begin{align}
    e_{ij} &= \text{LeakyReLU}\left(\mathbf{a}^T [\mathbf{W}h_i \| \mathbf{W}h_j]\right) \\
    \alpha_{ij} &= \frac{\exp(e_{ij})}{\sum_{k \in \mathcal{N}(i)} \exp(e_{ik})} \\
    h'_i &= \sigma\left(\sum_{j \in \mathcal{N}(i)} \alpha_{ij} \mathbf{W} h_j\right)
\end{align}

\begin{lstlisting}[caption={Custom GAT layer (src/models/gat\_model.py)}]
class GATLayer(nn.Module):
    """Single-head graph attention layer."""
    def __init__(self, in_features, out_features,
                 negative_slope=0.2, dropout=0.1):
        super().__init__()
        self.W = nn.Linear(in_features, out_features,
                           bias=False)
        self.a = nn.Parameter(
            torch.empty(2 * out_features))
        nn.init.xavier_uniform_(self.a.unsqueeze(0))
        self.leaky_relu = nn.LeakyReLU(negative_slope)
        self.dropout = nn.Dropout(dropout)

    def forward(self, x, edge_index, edge_weight=None):
        # x: [N, in_features]
        # edge_index: [2, E] directed edges
        N = x.size(0)
        Wh = self.W(x)                  # [N, out_feat]
        src, tgt = edge_index            # each [E]

        # Attention coefficients
        Wh_src = Wh[src]                 # [E, out_feat]
        Wh_tgt = Wh[tgt]                # [E, out_feat]
        cat = torch.cat([Wh_src, Wh_tgt], dim=1)
        e = self.leaky_relu(cat @ self.a)  # [E]

        # Sparse softmax per target node
        e_max = torch.full((N,), float("-inf"),
                           device=x.device)
        e_max.scatter_reduce_(0, tgt, e,
            reduce="amax", include_self=False)
        e_exp = torch.exp(e - e_max[tgt])
        e_sum = torch.zeros(N, device=x.device)
        e_sum.scatter_add_(0, tgt, e_exp)
        alpha = e_exp / (e_sum[tgt] + 1e-10)
        alpha = self.dropout(alpha)

        # Weighted aggregation
        messages = alpha.unsqueeze(1) * Wh_src
        out = torch.zeros(N, Wh.size(1),
                          device=x.device)
        out.scatter_add_(0,
            tgt.unsqueeze(1).expand_as(messages),
            messages)
        return out
\end{lstlisting}

\begin{lstlisting}[caption={Graph-enhanced model with residual GAT (src/models/gat\_model.py)}]
class GraphEnhancedModel(nn.Module):
    def __init__(self, num_features=10,
                 encoder_dim=256, gat_hidden=64,
                 gat_heads=4, num_classes=2):
        super().__init__()
        # Shared CNN-BiLSTM price encoder
        self.conv = nn.Sequential(
            nn.Conv1d(num_features, 64, 3, padding=1),
            nn.ReLU(), nn.BatchNorm1d(64),
            nn.Conv1d(64, 64, 3, padding=1),
            nn.ReLU(), nn.BatchNorm1d(64))
        self.lstm = nn.LSTM(64, 128, num_layers=2,
            batch_first=True, dropout=0.2,
            bidirectional=True)
        # Projection + two-layer GAT with residuals
        gat_in_dim = gat_hidden * gat_heads
        self.node_proj = nn.Linear(encoder_dim,
                                    gat_in_dim)
        self.gat1 = MultiHeadGAT(gat_in_dim,
            gat_hidden, gat_heads)
        self.gat_norm1 = nn.LayerNorm(gat_in_dim)
        self.gat2 = MultiHeadGAT(gat_in_dim,
            gat_hidden, gat_heads)
        self.gat_norm2 = nn.LayerNorm(gat_in_dim)
        self.head = nn.Sequential(
            nn.Linear(gat_in_dim, 128), nn.ReLU(),
            nn.Dropout(0.2), nn.Linear(128, 2))

    def forward(self, x, edge_index, edge_weight=None):
        # x: [N, T, F] for N companies
        h = self.encode_price(x)        # [N, 256]
        h = self.node_proj(h)           # [N, gat_dim]
        # GAT layer 1 + residual + LayerNorm
        h1 = F.elu(self.gat1(h, edge_index))
        h = self.gat_norm1(h + h1)
        # GAT layer 2 + residual + LayerNorm
        h2 = F.elu(self.gat2(h, edge_index))
        h = self.gat_norm2(h + h2)
        return self.head(h)             # [N, 2]
\end{lstlisting}

\textbf{Result:} GAT contributes $+2.6\%$ AUC over a no-GAT ablation, confirming that inter-company information propagation adds discriminative signal.

\subsection{Model 3: FinBERT Document Encoder}

\begin{lstlisting}[caption={FinBERT with attention pooling (src/models/document\_model.py)}]
class AttentionPooling(nn.Module):
    def __init__(self, hidden_size):
        super().__init__()
        self.score = nn.Linear(hidden_size, 1)

    def forward(self, sequence_output, attention_mask):
        scores = self.score(sequence_output).squeeze(-1)
        scores = scores.masked_fill(
            attention_mask == 0, float("-inf"))
        weights = torch.softmax(scores, dim=-1) \
            .unsqueeze(-1)
        return torch.sum(
            sequence_output * weights, dim=1)

class DocumentDirectionModel(nn.Module):
    def __init__(self, backbone="ProsusAI/finbert",
                 dropout=0.2):
        super().__init__()
        self.encoder = AutoModel.from_pretrained(
            backbone)
        hidden = self.encoder.config.hidden_size  # 768
        self.pool = AttentionPooling(hidden)
        self.head = nn.Sequential(
            nn.Dropout(dropout),
            nn.Linear(hidden, 256),
            nn.ReLU(),
            nn.Dropout(dropout),
            nn.Linear(256, 2),
        )

    def forward(self, input_ids, attention_mask):
        output = self.encoder(
            input_ids=input_ids,
            attention_mask=attention_mask)
        pooled = self.pool(
            output.last_hidden_state, attention_mask)
        return self.head(pooled)
\end{lstlisting}

\textbf{Why frozen backbone?} With only 227 training filings, fine-tuning FinBERT's 110M parameters would cause catastrophic overfitting. Only the attention pooling layer and head (small number of parameters) are trained.

\textbf{Output:} 768-dimensional embedding per (stock, year) pair.

\subsection{Model 4: Macro MLP Encoder}

\begin{lstlisting}[caption={Macro feature encoder (src/models/macro\_model.py)}]
class MacroStateModel(nn.Module):
    """3-layer MLP: 12 -> 64 -> 64 -> 32."""
    def __init__(self, input_dim=12, hidden_dim=64,
                 output_dim=32, dropout=0.2):
        super().__init__()
        self.encoder = nn.Sequential(
            nn.Linear(input_dim, hidden_dim),
            nn.LayerNorm(hidden_dim),
            nn.GELU(),
            nn.Dropout(dropout),
            nn.Linear(hidden_dim, hidden_dim),
            nn.LayerNorm(hidden_dim),
            nn.GELU(),
            nn.Dropout(dropout),
            nn.Linear(hidden_dim, output_dim),
        )

    def encode(self, x):
        return self.encoder(x)  # [B, 12] -> [B, 32]
\end{lstlisting}

\textbf{Output:} 32-dimensional embedding encoding the current macroeconomic regime.

\subsection{Fusion Model with Gated Attention}

\subsubsection{Architecture}

The fusion model combines all four modality embeddings through a learned sigmoid gating mechanism:

\begin{equation}
    g_i = \frac{\sigma(f_{\text{gate}}(\mathbf{z}))_i \cdot m_i}{\sum_{j=1}^{4} \sigma(f_{\text{gate}}(\mathbf{z}))_j \cdot m_j + \epsilon}
\end{equation}

where $m_i \in \{0,1\}$ is the availability mask, $\sigma$ is sigmoid, and $f_{\text{gate}}$ is a two-layer MLP.

\begin{lstlisting}[caption={Phase 14 Fusion Model with HAR-RV skip (src/models/fusion\_model.py)}]
class Phase14FusionModel(nn.Module):
    """Adds HAR-RV skip connection to gated fusion."""
    NUM_MODALITIES = 4

    def __init__(self, price_dim=256, har_rv_dim=3,
                 har_proj_dim=32, gat_dim=256,
                 doc_dim=768, macro_dim=32,
                 surprise_dim=5, proj_dim=128,
                 hidden_dim=256, dropout=0.3):
        super().__init__()
        # HAR-RV skip (direct, ungated)
        self.har_rv_skip = nn.Sequential(
            nn.Linear(har_rv_dim, har_proj_dim),
            nn.LayerNorm(har_proj_dim), nn.GELU())
        # Per-modality projectors
        def _proj(in_dim):
            return nn.Sequential(
                nn.Linear(in_dim, proj_dim),
                nn.LayerNorm(proj_dim),
                nn.GELU(),
                nn.Dropout(dropout * 0.5))
        self.price_proj = _proj(price_dim)
        self.gat_proj   = _proj(gat_dim)
        self.doc_proj   = _proj(doc_dim)
        self.macro_proj = _proj(macro_dim)
        # Gating (surprise-conditioned)
        gate_in = proj_dim * 4 + surprise_dim
        self.gate = nn.Sequential(
            nn.Linear(gate_in, hidden_dim), nn.GELU(),
            nn.Linear(hidden_dim, 4), nn.Sigmoid())
        # Shared trunk: gated + HAR-RV
        trunk_in = proj_dim * 4 + har_proj_dim
        self.trunk = nn.Sequential(
            nn.Linear(trunk_in, hidden_dim),
            nn.LayerNorm(hidden_dim), nn.GELU(),
            nn.Dropout(dropout),
            nn.Linear(hidden_dim, hidden_dim),
            nn.LayerNorm(hidden_dim), nn.GELU(),
            nn.Dropout(dropout * 0.5))
        # Volatility head (Softplus ensures > 0)
        self.volatility_head = nn.Sequential(
            nn.Linear(hidden_dim, hidden_dim // 2),
            nn.GELU(), nn.Dropout(dropout * 0.5),
            nn.Linear(hidden_dim//2, hidden_dim//4),
            nn.GELU(),
            nn.Linear(hidden_dim // 4, 1),
            nn.Softplus())
        # Direction head
        self.direction_head = nn.Sequential(
            nn.Linear(hidden_dim, hidden_dim // 4),
            nn.GELU(), nn.Linear(hidden_dim//4, 2))
\end{lstlisting}

\begin{lstlisting}[caption={Forward pass with gating and HAR-RV skip}]
    def forward(self, price_emb, har_rv_raw, gat_emb,
                doc_emb, macro_emb, surprise_feat,
                modality_mask):
        p = self.price_proj(price_emb)
        g = self.gat_proj(gat_emb)
        d = self.doc_proj(doc_emb)
        m = self.macro_proj(macro_emb)
        har = self.har_rv_skip(har_rv_raw)  # [B, 32]

        stacked = torch.stack([p, g, d, m], dim=1)
        stacked = stacked * modality_mask.unsqueeze(-1)
        flat = stacked.view(stacked.size(0), -1)
        gate_in = torch.cat([flat, surprise_feat], 1)
        gates = self.gate(gate_in)
        gates = gates * modality_mask
        gates = gates / (gates.sum(1, True) + 1e-8)

        weighted = stacked * gates.unsqueeze(-1)
        fused = weighted.view(weighted.size(0), -1)
        # Concatenate gated + HAR-RV skip
        trunk_in = torch.cat([fused, har], dim=-1)
        h = self.trunk(trunk_in)

        vol = self.volatility_head(h).squeeze(-1)
        dir_logits = self.direction_head(h)
        return {"volatility_pred": vol,
                "direction_logits": dir_logits,
                "gate_weights": gates}
\end{lstlisting}

\textbf{HAR-RV Skip Connection (Phase 14 key innovation):}
The CNN-BiLSTM compresses $[60, 21] \to 256$ dimensions, losing some autoregressive signal. The skip connection extracts \texttt{rv\_lag1d}, \texttt{rv\_lag5d}, \texttt{rv\_lag22d} from the last timestep and feeds them \textit{directly} into the fusion trunk, bypassing the bottleneck. This single addition improved volatility $R^2$ from 0.867 to 0.921.

\subsection{Loss Functions}

\subsubsection{QLIKE Loss for Volatility (Weight: 0.85)}

\begin{equation}
    \mathcal{L}_{\text{QLIKE}} = \frac{1}{N}\sum_{i=1}^N \left(\frac{\sigma^2_{\text{true},i}}{\hat{\sigma}^2_i} + \log\hat{\sigma}^2_i\right)
\end{equation}

\textbf{Why QLIKE over MSE?} When $\hat{\sigma} < \sigma_{\text{true}}$ (underprediction), the first term $\sigma^2_{\text{true}}/\hat{\sigma}^2$ explodes. Underpredicting risk is more dangerous than overpredicting it. QLIKE is also scale-invariant.

\begin{lstlisting}[caption={QLIKE loss (src/models/losses.py)}]
class QLIKELoss(nn.Module):
    """QLIKE (Quasi-Likelihood) loss for volatility
    forecasting (Patton, 2011)."""
    def __init__(self, eps=1e-6):
        super().__init__()
        self.eps = eps

    def forward(self, pred, target):
        valid = ((target > self.eps)
                 & (pred > self.eps)
                 & ~torch.isnan(target))
        if not valid.any():
            return torch.tensor(0.0, device=pred.device,
                                requires_grad=True)
        pred_var = pred[valid].pow(2).clamp(min=self.eps)
        true_var = target[valid].pow(2).clamp(
            min=self.eps)
        qlike = true_var / pred_var + torch.log(pred_var)
        return qlike.mean()
\end{lstlisting}

\subsubsection{ListNet Ranking Loss for Direction (Weight: 0.15)}

\begin{equation}
    \mathcal{L}_{\text{ListNet}} = \frac{1}{|D|}\sum_{d \in D} \text{CE}\!\left(\text{softmax}(\mathbf{y}_d \cdot \tau) \;\big\|\; \text{softmax}(\hat{\mathbf{y}}_d)\right)
\end{equation}

\textbf{Why ListNet over Cross-Entropy?} Cross-entropy evaluates each stock independently. ListNet groups all 30 stocks on the same trading date and optimizes their \textit{relative ranking}, canceling out market-wide movement (beta) and forcing the model to learn stock-specific signals.

\begin{lstlisting}[caption={ListNet ranking loss (src/models/losses.py)}]
def listnet_loss(scores, labels, dates,
                 temperature=5.0):
    """ListNet ranking loss.
    Groups stocks by trading date and computes
    cross-entropy between predicted and true rankings.
    """
    pred_probs = F.softmax(scores, dim=1)[:, 1]
    loss = torch.tensor(0.0, device=scores.device)
    n_dates = 0
    for date in dates.unique():
        mask = dates == date
        n_stocks = mask.sum()
        if n_stocks < 2:
            continue
        pred = pred_probs[mask]
        true = labels[mask].float()
        pred_dist = F.softmax(pred, dim=0)
        true_dist = F.softmax(true * temperature,
                              dim=0)
        ce = -(true_dist
               * torch.log(pred_dist + 1e-8)).sum()
        loss = loss + ce
        n_dates += 1
    return loss / max(n_dates, 1)
\end{lstlisting}

\subsubsection{Combined Loss}

\begin{equation}
    \mathcal{L} = 0.85 \times \mathcal{L}_{\text{QLIKE}} + 0.15 \times \mathcal{L}_{\text{ListNet}}
\end{equation}

\textbf{Why 0.85/0.15?} Three variants were tested: balanced (0.5/0.5), dir-assisted (0.7/0.3), and vol-primary (0.85/0.15). The balanced variant stopped training at epoch 10 because QLIKE and ListNet gradients pointed in opposing directions. The 0.85/0.15 split gave one objective a dominant voice, allowing the optimizer to converge (training continued to epoch 26).

\subsection{Training Configuration}

\begin{table}[H]
\centering
\caption{Hyperparameters for the fusion model}
\begin{tabular}{lll}
\toprule
\textbf{Hyperparameter} & \textbf{Value} & \textbf{Justification} \\
\midrule
Optimizer         & AdamW            & Correct L2 regularization (vs Adam) \\
Learning rate     & $5 \times 10^{-4}$ & Standard for Adam-family on financial data \\
Weight decay      & $10^{-3}$        & L2 regularization to prevent overfitting \\
Batch size        & 2048             & Large batch for stable ListNet gradient \\
LR Scheduler      & CosineAnnealing  & Smooth decay, avoids sharp LR drops \\
Max epochs        & 60               & With early stopping \\
Early stopping    & Patience = 7     & Stops if val loss flat for 7 epochs \\
Gradient clipping & max\_norm = 1.0  & Prevents gradient explosion \\
AMP               & fp16             & Halves memory, enables larger batches \\
\bottomrule
\end{tabular}
\end{table}

\begin{lstlisting}[caption={Training infrastructure (src/train/common.py)}]
@dataclass
class TrainingConfig:
    learning_rate: float = 1e-4
    weight_decay: float = 1e-4
    epochs: int = 20
    batch_size: int = 512
    patience: int = 5
    use_amp: bool = True
    gradient_accumulation_steps: int = 1

def create_optimizer(model, config):
    return torch.optim.AdamW(
        model.parameters(),
        lr=config.learning_rate,
        weight_decay=config.weight_decay)

class EarlyStopping:
    """Stop when monitored metric stops improving."""
    def __init__(self, patience=5, mode="min",
                 min_delta=1e-4):
        self.patience = patience
        self.mode = mode
        self.min_delta = min_delta
        self.counter = 0
        self.best_score = None
        self.should_stop = False

    def __call__(self, score):
        if self.best_score is None:
            self.best_score = score
            return False
        improved = (score < self.best_score
                    - self.min_delta
                    if self.mode == "min"
                    else score > self.best_score
                    + self.min_delta)
        if improved:
            self.best_score = score
            self.counter = 0
        else:
            self.counter += 1
            if self.counter >= self.patience:
                self.should_stop = True
        return self.should_stop
\end{lstlisting}

\subsection{Results}

\subsubsection{Volatility Forecasting Performance}

\begin{table}[H]
\centering
\caption{Volatility R\textsuperscript{2} progression across all project phases}
\begin{tabular}{lrr}
\toprule
\textbf{Phase / Model} & \textbf{Vol R\textsuperscript{2}} & \textbf{Dir AUC} \\
\midrule
Historical Average (naive baseline) & 0.348 & --- \\
V1 Fusion (Phase 6)                 & 0.415 & 0.516 \\
V2 + ListNet (Phase 10)             & 0.335 & 0.599 \\
QLIKE + 30 Stocks (Phase 12)        & 0.772 & 0.568 \\
HAR-RV Features (Phase 13)          & 0.867 & 0.585 \\
\textbf{HAR-RV Skip (Phase 14)}     & \textbf{0.921} & \textbf{0.591} \\
HAR-RV Benchmark                    & 0.947 & --- \\
\bottomrule
\end{tabular}
\end{table}

\begin{figure}[H]
\centering
\includegraphics[width=\textwidth]{figures/r2_progression.png}
\caption{Volatility R\textsuperscript{2} improvement across phases. Phase 14 closes 95.7\% of the gap between V2 baseline and the HAR-RV benchmark.}
\end{figure}

\subsubsection{Learned Gate Weights}

\begin{figure}[H]
\centering
\includegraphics[width=0.7\textwidth]{figures/gate_weights.png}
\caption{Fusion gate weights in the final model. GAT (41.2\%) is dominant, confirming that cross-stock relationships are the most informative signal once data leakage is removed.}
\end{figure}

\subsubsection{Backtest Results (Direction Strategy)}

A weekly long-short strategy using the direction head:
\begin{itemize}[noitemsep]
    \item \textbf{Setup:} Top-2 long / Bottom-2 short by predicted UP probability, weekly rebalancing
    \item \textbf{Sharpe at 0bp:} 0.697 \quad \textbf{At 5bp:} 0.620 \quad \textbf{At 10bp:} 0.543
    \item \textbf{Break-even:} $\approx$20bp --- strategy remains viable at realistic transaction costs
\end{itemize}

%══════════════════════════════════════════════════════════════════════════════
\section{Summary \& Next Steps}
%══════════════════════════════════════════════════════════════════════════════

\subsection{Phase 3 Deliverables Completed}
\begin{itemize}[noitemsep]
    \item[$\checkmark$] Chronological train/validation/test split (65.8/10.9/23.3\%)
    \item[$\checkmark$] Z-score normalization (fit on training set only)
    \item[$\checkmark$] 21-feature engineering pipeline including HAR-RV autoregressive features
    \item[$\checkmark$] FinBERT chunking pipeline for 227 SEC filings
    \item[$\checkmark$] Macroeconomic feature extraction with 1-day lag
    \item[$\checkmark$] 30-node knowledge graph construction
    \item[$\checkmark$] Availability masking for missing modalities
    \item[$\checkmark$] Two data leakage audits and corrections
\end{itemize}

\subsection{Phase 4 Deliverables Completed}
\begin{itemize}[noitemsep]
    \item[$\checkmark$] CNN-BiLSTM price encoder (256-d output)
    \item[$\checkmark$] Custom GAT in pure PyTorch (256-d output, +2.6\% AUC over ablation)
    \item[$\checkmark$] Frozen FinBERT with attention pooling (768-d output)
    \item[$\checkmark$] Macro MLP encoder (32-d output)
    \item[$\checkmark$] Gated multimodal fusion with availability masking
    \item[$\checkmark$] HAR-RV skip connection ($R^2$: 0.867 $\to$ 0.921)
    \item[$\checkmark$] QLIKE loss for volatility + ListNet ranking loss for direction
    \item[$\checkmark$] AMP fp16 training with gradient clipping
\end{itemize}

\subsection{Key Lessons Learned}
\begin{enumerate}
    \item \textbf{Data integrity is paramount.} Two separate leakages demonstrated how easily an ML pipeline can appear to work when a feature encodes the target. Auditing inputs when results seem too good is essential.
    \item \textbf{Task choice matters.} Volatility is more predictable than direction on efficient markets. Switching the primary objective from direction to volatility increased $R^2$ from 0.335 to 0.921.
    \item \textbf{Skip connections preserve information.} The HAR-RV skip connection directly addressed the CNN-BiLSTM compression bottleneck, providing the single largest improvement ($+6.3\%$ $R^2$).
    \item \textbf{Loss function design is architecture.} QLIKE's asymmetric penalty and scale-invariance, and ListNet's beta-canceling ranking formulation, contributed as much as architectural changes.
\end{enumerate}

\vspace{1cm}
\noindent\textbf{GitHub Repository:} \url{https://github.com/imrishi007/financial-document-analysis}


\end{document}
